\documentclass[11pt]{amsart}

\usepackage{amsmath,amssymb,amsthm}
\usepackage{mathtools}
\usepackage{enumitem}
\usepackage{booktabs}
\usepackage{array}
\usepackage{hyperref}
\usepackage{cleveref}

% Theorem environments
\newtheorem{theorem}{Theorem}[section]
\newtheorem{lemma}[theorem]{Lemma}
\newtheorem{proposition}[theorem]{Proposition}
\newtheorem{corollary}[theorem]{Corollary}
\newtheorem{conjecture}[theorem]{Conjecture}

\theoremstyle{definition}
\newtheorem{definition}[theorem]{Definition}
\newtheorem{example}[theorem]{Example}
\newtheorem{observation}[theorem]{Observation}

\theoremstyle{remark}
\newtheorem{remark}[theorem]{Remark}

% Notation
\newcommand{\R}{\mathbb{R}}
\newcommand{\Q}{\mathbb{Q}}
\newcommand{\Z}{\mathbb{Z}}
\newcommand{\N}{\mathbb{N}}
\newcommand{\vphi}{\varphi}
\newcommand{\Spec}{\operatorname{Spec}}

\title[The Tihany Conjecture: Spectral Resolutions]{%
The Erd\H{o}s--Lov\'asz Tihany Conjecture:\\
Spectral Resolutions}

\author{V.\,F.\,S.~Santos}
\date{February 2026}

\subjclass[2020]{05C15, 05C50, 05C31, 11R04}

\keywords{Tihany conjecture, Mycielski graphs, spectral graph theory, 
golden ratio, eigenvector partitions, chromatic number}

\begin{document}

\begin{abstract}
We apply golden resolvent theory~\cite{santos-grt-2026} to the 
Erd\H{o}s--Lov\'asz Tihany Conjecture (1968): for every graph~$G$ 
with $\chi(G) = s + t - 1 > \omega(G)$, there exists a partition 
$V(G) = S \sqcup T$ with $\chi(G[S]) \ge s$ and $\chi(G[T]) \ge t$.

We prove the conjecture for the entire Mycielski family: for every 
$k \ge 3$ and every pair $(s,t)$ with $s + t = k + 1$ and $s, t \ge 2$, 
the Mycielski graph~$M_k$ admits a valid Tihany partition.  The proof 
uses the canonical embedding $M_s \hookrightarrow M_k$ and a 
Mycielski-complement induction.  We extend this to all generalised 
Mycielskian chains via an abstract $(P1)$--$(P3)$ operator framework.

Beyond Mycielski graphs, we formulate two conjectures that each imply 
the full Tihany Conjecture.  The \emph{Spectral Tihany Conjecture} 
asserts that some eigenvector threshold cut of~$G$ itself yields a 
valid partition; we verify it on $161$ Tihany pairs with zero failures.  
The \emph{Mycielskian Spectral Pullback Conjecture} asserts that 
embedding~$G$ into its Mycielskian~$M(G)$ and sweeping with an 
eigenvector of~$M(G)$---whose spectrum is controlled by the secular 
equation of~\cite{santos-grt-2026}---always produces a valid partition 
when restricted to~$G$; we verify it on $94$ graphs with zero failures.

The spectral partition mechanisms---cooperative modes for 
regular graphs, resolvent centrality for irregular ones---follow from 
the complementarity principle of golden resolvent 
theory~\cite{santos-grt-2026}.
\end{abstract}

\maketitle


%% ===================================================================
\section{Introduction}\label{sec:intro}
%% ===================================================================

\subsection{The conjecture}

\begin{conjecture}[Erd\H{o}s--Lov\'asz, 1968]\label{conj:tihany}
For every graph~$G$ and integers $s, t \ge 2$ with 
$\chi(G) = s + t - 1 > \omega(G)$, there exists a partition 
$V(G) = S \sqcup T$ with $\chi(G[S]) \ge s$ and $\chi(G[T]) \ge t$.
\end{conjecture}

The conjecture has been open since 1968~\cite{erdos-lovasz-1975}.  
Known cases include small pairs with $s + t \le 6$~\cite{stiebitz-1987, 
sachs-1993}, line graphs~\cite{kostochka-stiebitz-2003}, quasi-line 
graphs~\cite{chudnovsky-fradkin-2008}, graphs with 
$\alpha(G) = 2$~\cite{balogh-etal-2009}, and the pair $(3, k{-}2)$ 
on Mycielski graphs for all $k \ge 5$~\cite{santos-golden-2026}.  
The general conjecture remains wide open.

\subsection{Main results}

The contributions of this paper are:
\begin{enumerate}[label=(\arabic*)]
\item \textbf{Tihany for all Mycielski pairs} (\Cref{thm:tihany-mycielski}): 
  for every $k \ge 3$ and every valid pair $(s,t)$, $M_k$ admits a 
  Tihany partition.  This is a complete resolution of the conjecture 
  for the Mycielski family.

\item \textbf{Extension to generalised Mycielskians} 
  (\Cref{prop:generalised}): any chain built by an operator satisfying 
  chromatic increment, induced embedding, and bounded complement 
  satisfies all Tihany pairs.

\item \textbf{The Spectral Tihany Conjecture} (\Cref{conj:spectral-tihany}): 
  for every graph with $\chi > \omega$ and every Tihany pair, some 
  eigenvector threshold cut produces a valid partition.  
  Verified: 161/161 pairs.

\item \textbf{The Mycielskian Spectral Pullback Conjecture} 
  (\Cref{conj:pullback}): for every graph with $\chi > \omega$, some 
  eigenvector of $M(G)$ produces a valid Tihany partition of~$G$ via 
  pullback sweep.  Verified: 94/94 graphs.

\item \textbf{Spectral shielding} (\Cref{prop:shielding}): the 
  Hoffman bound, Lov\'asz theta, and fractional chromatic number are 
  all permanently below~4 on Mycielski graphs---yet eigenvector 
  partitions succeed.
\end{enumerate}

\subsection{Relation to companion papers}

The spectral machinery used here---the golden resolvent factorisation, 
the secular equation, the complementarity principle, the Positive 
Boost Inequality, and the Galois orbit framework---is developed in 
the companion paper~\cite{santos-grt-2026} (\emph{Golden Resolvent 
Theory}) as pure mathematics applicable to arbitrary real symmetric 
matrices.  The present paper applies that theory to graph chromatic 
partitions.

The published paper~\cite{santos-golden-2026} proved the Tihany 
conjecture for the single family $(3, k{-}2)$ on Mycielski graphs 
via the $C_5$-peeling method and the Golden Sub-Induction.  
\Cref{thm:tihany-mycielski} of the present paper supersedes that 
result, resolving all $k{-}2$ Tihany pairs simultaneously.

\subsection{Organisation}

\Cref{sec:spectral-tihany} introduces eigenvector threshold cuts 
and states the Spectral Tihany Conjecture.  \Cref{sec:bounds} 
discusses spectral lower bounds and the shielding phenomenon.  
\Cref{sec:mycielski} proves Tihany for all Mycielski pairs.  
\Cref{sec:generalised} extends to generalised Mycielskians.  
\Cref{sec:pullback} develops the Mycielskian Spectral Pullback.  
\Cref{sec:evidence} presents the computational evidence.  
\Cref{sec:open} lists open problems.


%% ===================================================================
\section{Eigenvector Threshold Cuts and the Spectral Tihany Conjecture}
\label{sec:spectral-tihany}
%% ===================================================================

Let $G = (V, E)$ be a simple graph on $n$~vertices with adjacency 
matrix $A \in \Z^{n \times n}$.  Since $A$ is real symmetric, it has 
$n$ real eigenvalues $\lambda_1 \ge \cdots \ge \lambda_n$ with an 
orthonormal eigenbasis $v_1, \ldots, v_n$.

\begin{definition}[Threshold cut]\label{def:threshold}
For an eigenvector $v \in \R^n$ and threshold $\tau \in \R$, the 
\emph{threshold cut} is the partition 
$S_\tau = \{i \in V : v_i \ge \tau\}$, $T_\tau = V \setminus S_\tau$.
\end{definition}

Since $S_\tau$ changes only when $\tau$ crosses an eigenvector component, 
there are at most $n-1$ distinct nontrivial partitions per eigenvector, 
and at most $n(n-1)$ across all eigenvectors.

\begin{conjecture}[Spectral Tihany Conjecture]\label{conj:spectral-tihany}
For every graph~$G$ with $\chi(G) > \omega(G)$ and every valid Tihany 
pair $(s,t)$ with $s + t = \chi(G) + 1$ and $s, t \ge 2$, there 
exists an eigenvector~$v$ of $A(G)$ and a threshold~$\tau$ such that 
$\chi(G[S_\tau]) \ge s$ and $\chi(G[T_\tau]) \ge t$.
\end{conjecture}

\begin{proposition}\label{prop:spectral-implies-tihany}
\Cref{conj:spectral-tihany} implies the Erd\H{o}s--Lov\'asz Tihany 
Conjecture (\Cref{conj:tihany}).
\end{proposition}

\begin{proof}
The threshold cut $(S_\tau, T_\tau)$ is a vertex partition satisfying 
the required chromatic lower bounds.
\end{proof}

\begin{remark}[Finite optimisation]\label{rmk:finite-opt}
\Cref{conj:spectral-tihany} reduces the Tihany problem to a finite 
search: for a given graph and pair, check all $O(n^2)$ candidate 
partitions.  All our computational results are exact verifications 
of this type.
\end{remark}


%% ===================================================================
\section{Spectral Lower Bounds and the Shielding Phenomenon}
\label{sec:bounds}
%% ===================================================================

\subsection{Hoffman bound}

\begin{definition}[Hoffman bound~\cite{hoffman-1970}]\label{def:hoffman}
For a graph~$H$ with $\lambda_{\min}(H) < 0$:
\[
  L_{\mathrm{Hof}}(H) = 1 - \frac{\lambda_{\max}(H)}{\lambda_{\min}(H)} 
  \;\le\; \chi(H).
\]
\end{definition}

\begin{definition}[Hoffman-certified success]\label{def:hoffman-cert}
An eigenvector $v$ of $A(G)$ \emph{Hoffman-certifies} the pair $(s,t)$ 
if some threshold cut produces a partition with 
$L_{\mathrm{Hof}}(G[S]) \ge s$ and $L_{\mathrm{Hof}}(G[T]) \ge t$.
\end{definition}

\subsection{The Hoffman certification gap}

\begin{definition}[Hoffman gap]\label{def:hoffman-gap}
$\Delta_{\mathrm{Hof}}(G) = \chi(G) - L_{\mathrm{Hof}}(G)$.
\end{definition}

For Mycielski graphs, the gap is dramatic: $\Delta_{\mathrm{Hof}}(M_6) = 3.46$.

\subsection{Spectral shielding}

\begin{proposition}[Spectral shielding]\label{prop:shielding}
The three standard spectral and convex-relaxation bounds---the 
Hoffman bound $L_{\mathrm{Hof}}$, the Lov\'asz theta 
function~$\bar{\vartheta}$, and the fractional chromatic 
number~$\chi_f$---satisfy
\[
  L_{\mathrm{Hof}}(M_k) \;\le\; \bar{\vartheta}(M_k) \;\le\; 
  \chi_f(M_k) \;<\; 4 \quad \text{for every } k \ge 3.
\]
In particular, none of these bounds can certify $\chi \ge 4$ for 
any Mycielski graph.
\end{proposition}

\begin{proof}
The Larsen--Propp--Ullman recurrence~\cite{larsen-propp-ullman-1995} 
gives $\chi_f(M(G)) = \chi_f(G) + 1/\alpha(G)$.  Since 
$\alpha(M_k) = n_{k-1}$, the fractional chromatic number converges:
\[
  \chi_f(M_\infty) = \tfrac{5}{2} + \sum_{j=3}^\infty 
  \frac{1}{\alpha(M_j)} \approx 3.377 < 4.
\]
By the Lov\'asz sandwich theorem, 
$L_{\mathrm{Hof}} \le \bar{\vartheta} \le \chi_f$, so all three 
are permanently below~4 as $\chi(M_k) = k \to \infty$.
\end{proof}

\begin{remark}
The shielding is caused by the abundance of independent shadow vertices.  
Yet eigenvector threshold cuts produce correct Tihany partitions in 
every tested case (\Cref{sec:evidence}).  The partition information 
resides in the \emph{geometry} of the eigenvector---the spatial 
pattern of signs and magnitudes---rather than in any eigenvalue ratio.
\end{remark}


%% ===================================================================
\section{Tihany for All Mycielski Pairs}\label{sec:mycielski}
%% ===================================================================

\subsection{The Mycielski construction}

The Mycielskian~\cite{mycielski-1955} of a graph~$G$ on $n$~vertices 
is $M(G)$ on $2n+1$ vertices: originals $V_0$, shadows $V_1$ 
(where $u_i' \sim N_G(v_i)$), and an apex~$w$ adjacent to all shadows.  
Key properties: $\chi(M(G)) = \chi(G) + 1$ and $\omega(M(G)) = \omega(G)$.

\subsection{Canonical embedding and complement}

\begin{definition}[Canonical chain embedding]\label{def:canonical}
Define $M_2 = K_2$ and $M_{k+1} = M(M_k)$.  The original block of 
$M_{k+1}$ is isomorphic to $M_k$, giving a canonical inclusion 
$\iota_{\ell,k}\colon V(M_\ell) \hookrightarrow V(M_k)$ for 
$2 \le \ell \le k$.
\end{definition}

\begin{definition}[Mycielski complement]\label{def:complement}
For $2 \le s \le k-1$:
\[
  R_{k,s} := M_k\bigl[V(M_k) \setminus \iota_{s,k}(V(M_s))\bigr].
\]
\end{definition}

\subsection{Main theorem}

\begin{lemma}[Containment]\label{lem:containment}
For every $k \ge s + 2$, the complement $R_{k,s}$ contains $M(R_{k-1,s})$ 
as a subgraph.
\end{lemma}

\begin{proof}
Write $M_k = M(M_{k-1})$ with vertex set $V_0 \sqcup V_1 \sqcup \{w\}$.  
Since $M_s$ embeds in $V_0 \cong M_{k-1}$, the complement decomposes as 
$R_{k,s} = R_{k-1,s} \cup V_1 \cup \{w\}$.  We exhibit $M(R_{k-1,s})$ 
inside $R_{k,s}$:
\begin{itemize}
\item Originals of $M(R_{k-1,s}) \to R_{k-1,s} \subset R_{k,s}$.
\item Shadow of $i \in R_{k-1,s} \to u_i \in V_1$; adjacencies preserved 
  since $u_i \sim N_{M_{k-1}}(i) \cap R_{k-1,s} = N_{R_{k-1,s}}(i)$.
\item Apex of $M(R_{k-1,s}) \to w$; adjacent to all shadows.
\end{itemize}
The map is injective and preserves all adjacencies.
\end{proof}

\begin{lemma}[Chromatic lower bound]\label{lem:chromatic}
For every $k \ge s + 1$ with $s \ge 2$: $\chi(R_{k,s}) \ge k - s + 1$.
\end{lemma}

\begin{proof}
Induction on~$k$.  Base case ($k = s+1$): $R_{s+1,s}$ is the star 
$K_{1,|V(M_s)|}$, with $\chi = 2 = (s+1) - s + 1$.  Inductive step 
($k \ge s+2$): by \Cref{lem:containment}, $R_{k,s} \supseteq M(R_{k-1,s})$.  
By the Mycielski chromatic increment, 
$\chi(R_{k,s}) \ge \chi(M(R_{k-1,s})) = \chi(R_{k-1,s}) + 1 \ge (k-s) + 1$.
\end{proof}

\begin{theorem}[Tihany for all Mycielski pairs]\label{thm:tihany-mycielski}
For every $k \ge 3$ and every pair $(s,t)$ with $s + t = k + 1$ and 
$s, t \ge 2$, the Mycielski graph~$M_k$ admits a partition 
$V(M_k) = S \sqcup T$ with $\chi(M_k[S]) \ge s$ and $\chi(M_k[T]) \ge t$.
\end{theorem}

\begin{proof}
Set $S = \iota_{s,k}(V(M_s))$ and $T = V(M_k) \setminus S$.  
Since $M_s$ is an induced subgraph, $\chi(M_k[S]) = \chi(M_s) = s$.  
By \Cref{lem:chromatic}, $\chi(M_k[T]) = \chi(R_{k,s}) \ge k - s + 1 = t$.  
Since $\chi(M_k) = k = s + t - 1$ and $\omega(M_k) = 2 < k$ for 
$k \ge 3$, this is a valid Tihany partition.
\end{proof}

\begin{remark}[Discovery through spectral theory]\label{rmk:discovery}
\Cref{thm:tihany-mycielski} is purely combinatorial, but it was 
\emph{discovered} through the secular equation 
of~\cite{santos-grt-2026}.  The golden-flip eigenvector at eigenvalue 
$(-1)^{k-2}\vphi^{-(k-2)}$---whose block structure 
$(v, -\vphi v, 0)$ places originals and shadows in opposite 
partition halves---was verified computationally to achieve every 
Tihany pair for $M_4$ through $M_6$.  Analysing the block structure 
of these partitions revealed the containment 
$R_{k,s} \supseteq M(R_{k-1,s})$, which is the engine of the proof.
\end{remark}

\begin{remark}[Comparison with $C_5$-peeling]\label{rmk:c5-peeling}
The published paper~\cite{santos-golden-2026} proved the single family 
$(3, k{-}2)$ via $C_5$-peeling and sub-induction.  
\Cref{thm:tihany-mycielski} supersedes it: it resolves all $k-2$ 
pairs simultaneously, the proof is shorter (two lemmas plus the 
Mycielski increment), and no eigenvector computation is required.  
The two approaches are complementary: $C_5$-peeling provides a 
spectral partition (useful for generalisations to non-Mycielski graphs), 
while \Cref{thm:tihany-mycielski} provides a structural partition 
(optimal for the Mycielski family).
\end{remark}


%% ===================================================================
\section{Extension to Generalised Mycielskians}\label{sec:generalised}
%% ===================================================================

The proof of \Cref{thm:tihany-mycielski} uses exactly three properties 
of the Mycielski operator:

\begin{enumerate}[label=(P\arabic*)]
\item \textbf{Chromatic increment}: $\chi(M(G)) = \chi(G) + 1$ for 
  every graph~$G$ with an edge.
\item \textbf{Induced embedding}: $G$ embeds as an induced subgraph 
  (the original block) of $M(G)$.
\item \textbf{Bounded complement}: $\chi(M(G) \setminus G) = 2$.
\end{enumerate}

\begin{proposition}[Tihany for generalised chains]\label{prop:generalised}
Let $p \ge 1$, let $G_0$ be any graph with $\chi(G_0) \ge 2$ and 
$\omega(G_0) < \chi(G_0)$, and define $G_k = M^{(p)}(G_{k-1})$ 
for $k \ge 1$.  Then every $G_k$ satisfies the Tihany Conjecture 
for all pairs.
\end{proposition}

\begin{proof}
It suffices to verify (P1)--(P3) for $F = M^{(p)}$.

\emph{(P1)}: The standard Mycielski argument adapts to $p$~shadow 
layers; see~\cite{lin-wu-lam-gu-2006}.

\emph{(P2)}: Layer-0 vertices form an induced copy of~$G$.

\emph{(P3)}: $M^{(p)}(G) \setminus G$ consists of layers 
$V_1, \ldots, V_p$ (independent sets) with inter-layer edges, plus 
the apex.  A 2-colouring: odd layers and apex get colour~1, even layers 
get colour~2.

With (P1)--(P3) established, \Cref{lem:containment} and 
\Cref{lem:chromatic} carry over with $M$ replaced by $M^{(p)}$.
\end{proof}

\begin{remark}[Scope and limitations]\label{rmk:scope}
The (P1)--(P3) framework captures ``iterated Mycielskian-type'' 
constructions.  It does not reach:
\begin{itemize}
\item \emph{Kneser graphs}: the natural embedding 
  $K(n{-}1,r) \hookrightarrow K(n,r)$ has complement with 
  $\chi = 1$ (pairwise intersecting sets), so (P3) fails.
\item \emph{General triangle-free graphs}: non-constructive families 
  lack the recursive block structure for~(P2).
\end{itemize}
The Spectral Tihany Conjecture (\Cref{conj:spectral-tihany}) targets 
these families via eigenvector threshold cuts.
\end{remark}


%% ===================================================================
\section{The Mycielskian Spectral Pullback}\label{sec:pullback}
%% ===================================================================

A third approach bypasses eigenvectors of~$G$ and instead embeds 
$G$ into $M(G)$, exploiting the structured spectrum of $M(G)$ 
via the secular equation of~\cite{santos-grt-2026}.

\begin{definition}[Pullback sweep]\label{def:pullback}
Let $G$ have vertex set $V_0 = \{0,\ldots,n-1\}$, and let $M(G)$ 
have vertex set $V_0 \cup V_1 \cup \{w\}$.  For an eigenvector~$u$ 
of $A(M(G))$, order all $2n+1$ vertices by decreasing component value.  
At step~$k$, restrict the partition $(S_k, T_k)$ to the original block:
\[
  S_k^G = S_k \cap V_0, \qquad T_k^G = T_k \cap V_0.
\]
\end{definition}

\begin{conjecture}[Mycielskian Spectral Pullback]
\label{conj:pullback}
For every graph~$G$ with $\chi(G) > \omega(G)$ and every Tihany pair 
$(s,t)$, there exists an eigenvector~$u$ of $A(M(G))$ and a step~$k$ 
in the pullback sweep such that 
$\chi(G[S_k^G]) \ge s$ and $\chi(G[T_k^G]) \ge t$.
\end{conjecture}

\begin{proposition}\label{prop:pullback-implies}
\Cref{conj:pullback} implies the Erd\H{o}s--Lov\'asz Tihany 
Conjecture.
\end{proposition}

\subsection{The spectral mechanism}

By the secular equation~\cite[Theorem~4.1]{santos-grt-2026}, the 
spectrum of $M(G)$ decomposes into:

\begin{enumerate}[label=(\alph*)]
\item \textbf{Transparent eigenvalues} (from eigenspaces of~$G$ 
  orthogonal to~$\mathbf{1}$): produce golden pairs $\mu\vphi$ and 
  $-\mu/\vphi$ with eigenvectors $(v, v/\vphi, 0)$ 
  (\emph{Golden-Boost}) and $(v, -\vphi v, 0)$ (\emph{Golden-Flip}).

\item \textbf{Coupled eigenvalues} (from the secular equation): 
  eigenvectors with nonzero apex component; restriction to~$V_0$ is 
  proportional to resolvent centrality $(\lambda I - A)^{-1}\mathbf{1}$.
\end{enumerate}

By the complementarity principle~\cite[Propositions~4.5--4.6]{santos-grt-2026}:
\begin{itemize}
\item For \emph{regular} graphs: coupled eigenvectors are constant 
  on originals (useless); transparent Golden-Flip vectors produce 
  effective partitions.
\item For \emph{irregular} graphs: coupled eigenvectors sort originals 
  by density (effective); transparent vectors may be noisy.
\end{itemize}

\subsection{Cooperative-mode analysis}

Tracing each transparent eigenvector of $M(G)$ back to its parent 
eigenvalue~$\mu$ in~$G$ reveals a striking pattern.

\begin{definition}\label{def:cooperative}
An eigenvalue $\mu > 0$ of~$G$ is \emph{cooperative} (adjacent vertices 
receive same-sign components); $\mu < 0$ is \emph{frustrated} 
(alternating signs).
\end{definition}

\begin{center}
\begin{tabular}{lcc}
\toprule
Parent eigenvalue & Gap-free rate & Certifies Tihany \\
\midrule
Cooperative ($\mu > 0$), boost branch & 96\% & 96\% \\
Cooperative ($\mu > 0$), flip branch & 97\% & 97\% \\
Frustrated ($\mu < 0$), boost branch & 53\% & 53\% \\
Frustrated ($\mu < 0$), flip branch & 63\% & 63\% \\
\bottomrule
\end{tabular}
\end{center}

Cooperative modes of~$G$ produce nearly perfect partitions regardless 
of the secular branch.

\begin{conjecture}[Cooperative pullback for regular graphs]
\label{conj:cooperative}
For every $k$-regular graph~$G$ with $\chi(G) > \omega(G)$ and every 
Tihany pair $(s,t)$, there exists a cooperative transparent eigenvector 
of $M(G)$ whose pullback sweep certifies $(s,t)$.
\end{conjecture}


%% ===================================================================
\section{Computational Evidence}\label{sec:evidence}
%% ===================================================================

\subsection{Spectral Tihany Conjecture}

\begin{center}
\begin{tabular}{lrc}
\toprule
Suite & Pairs & Result \\
\midrule
Structured (20 graphs, 8 families) & 43 & 43/43 \\
Random (perturbed + Erd\H{o}s--R\'enyi) & 118 & 118/118 \\
\midrule
\textbf{Total} & \textbf{161} & \textbf{161/161} \\
\bottomrule
\end{tabular}
\end{center}

The structured suite spans odd cycles, odd antiholes, Mycielski chains, 
Kneser graphs, Platonic graphs, and Cartesian products; 
chromatic numbers range from 3 to~11, graph sizes up to $n = 84$.  
The Hoffman bound certifies 33 of 43 structured pairs (77\%); the 
remaining 10 are verified by exact chromatic number computation.

\subsection{Mycielskian Spectral Pullback}

\begin{center}
\begin{tabular}{lrc}
\toprule
Suite & Graphs & Result \\
\midrule
Named graphs (14 families) & 14 & 14/14 \\
Random Erd\H{o}s--R\'enyi ($\chi > \omega$) & 50 & 50/50 \\
Random triangle-free ($\chi > \omega$) & 30 & 30/30 \\
\midrule
\textbf{Total} & \textbf{94} & \textbf{94/94} \\
\bottomrule
\end{tabular}
\end{center}

\subsection{One-step gap analysis}

\begin{center}
\begin{tabular}{lccc}
\toprule
Eigenvector class & Tested & Gap-free & Certifies \\
\midrule
Golden-Flip & 43 & 38 (88\%) & 38 (88\%) \\
Golden-Boost & 43 & 35 (81\%) & 35 (81\%) \\
Mixed transparent & 19 & 17 (90\%) & 17 (90\%) \\
Coupled & 26 & 15 (58\%) & 15 (58\%) \\
\midrule
All transparent & 105 & 90 (86\%) & 90 (86\%) \\
\bottomrule
\end{tabular}
\end{center}

Transparent eigenvectors have a gap-free rate of 86\% individually, 
far exceeding the coupled rate of 58\%.  Every tested graph has at 
least one gap-free eigenvector.

\subsection{Stress test on larger ensembles}

\begin{center}
\begin{tabular}{lccc}
\toprule
Parent eigenvalue class & Tested & Gap-free & Rate \\
\midrule
Cooperative ($\mu > 0$, transparent) & 241 & 236 & 98\% \\
Frustrated ($\mu < 0$, transparent) & 1111 & 1024 & 92\% \\
Coupled (non-transparent) & 970 & 890 & 92\% \\
\bottomrule
\end{tabular}
\end{center}

The cooperative mode's near-perfect rate (98\%) supports a proof 
strategy for \Cref{conj:pullback}: prove the cooperative case for 
regular graphs (where positive transparent eigenvalues are guaranteed), 
and show that the coupled Perron vector handles the irregular complement.


%% ===================================================================
\section{Open Problems}\label{sec:open}
%% ===================================================================

\subsubsection*{Toward resolving the Tihany Conjecture.}

\begin{enumerate}
\item \textbf{Prove the Mycielskian Spectral Pullback} 
  (\Cref{conj:pullback}).  The proof splits naturally:
  \begin{enumerate}[label=(\alph*)]
  \item \emph{Regular graphs} (\Cref{conj:cooperative}): show that 
    cooperative transparent eigenvectors force the pullback sweep to 
    cross the $(s,t)$ threshold without a one-step gap.
  \item \emph{Irregular graphs}: show that the coupled Perron vector 
    of $M(G)$---sorting originals by resolvent centrality---always 
    admits a Tihany partition.
  \end{enumerate}

\item \textbf{Prove the Spectral Tihany Conjecture} 
  (\Cref{conj:spectral-tihany}) for any nontrivial infinite family.  
  The most accessible targets: odd antiholes $\bar{C}_{2m+1}$ and 
  Kneser graphs $K(n,k)$.

\item \textbf{Close the one-step gap.}  Prove a pair-indexed pivot 
  exclusion lemma: the vertex that raises $\chi(S)$ to~$s$ cannot 
  simultaneously drop $\chi(T)$ below~$t$ along some eigenvector ordering.

\item \textbf{Close the Hoffman gap.}  Find a spectral bound certifying 
  the 10~beyond-Hoffman cases, or prove that eigenvector geometry forces 
  valid partitions without any eigenvalue bound.

\item \textbf{Extend computational tests.}  Test $M_7$ ($n = 95$, 
  $\chi = 7$) and larger random graphs ($n \ge 50$).
\end{enumerate}


%% ===================================================================
%% REFERENCES
%% ===================================================================

\begin{thebibliography}{99}

\bibitem{balogh-etal-2009}
J.~Balogh, A.\,V.~Kostochka, N.~Prince, and M.~Stiebitz,
\emph{The Erd\H{o}s--Lov\'asz Tihany conjecture for quasi-line graphs},
Discrete Math.\ \textbf{309} (2009), no.~12, 3985--3991.

\bibitem{chudnovsky-fradkin-2008}
M.~Chudnovsky and A.~Fradkin,
\emph{An approximate version of a conjecture of Aharoni and Berger},
Preprint, 2008.

\bibitem{cvetkovic-rowlinson-simic-2010}
D.~Cvetkovi\'c, P.~Rowlinson, and S.~Simi\'c,
\emph{An Introduction to the Theory of Graph Spectra},
London Math.\ Soc.\ Student Texts, no.~75,
Cambridge Univ.\ Press, 2010.

\bibitem{erdos-lovasz-1975}
P.~Erd\H{o}s and L.~Lov\'asz,
\emph{Problems and results on chromatic numbers in finite and 
infinite graphs},
Recent Advances in Graph Theory, 1975, pp.~137--148.

\bibitem{hoffman-1970}
A.\,J.~Hoffman,
\emph{On eigenvalues and colorings of graphs},
Graph Theory and its Applications, Academic Press, 1970, pp.~79--91.

\bibitem{katz-1953}
L.~Katz,
\emph{A new status index derived from sociometric analysis},
Psychometrika \textbf{18} (1953), no.~1, 39--43.

\bibitem{kostochka-stiebitz-2003}
A.\,V.~Kostochka and M.~Stiebitz,
\emph{Colour-critical graphs with few edges},
Discrete Math.\ \textbf{273} (2003), 255--259.

\bibitem{larsen-propp-ullman-1995}
M.~Larsen, J.~Propp, and D.~Ullman,
\emph{The fractional chromatic number of Mycielski's graphs},
J.~Graph Theory \textbf{19} (1995), no.~3, 411--416.

\bibitem{lin-wu-lam-gu-2006}
W.~Lin, J.~Wu, P.\,C.\,B.~Lam, and G.~Gu,
\emph{Several parameters of generalized Mycielskians},
Discrete Appl.\ Math.\ \textbf{154} (2006), 1173--1182.

\bibitem{lovasz-1979}
L.~Lov\'asz,
\emph{On the Shannon capacity of a graph},
IEEE Trans.\ Inform.\ Theory \textbf{25} (1979), no.~1, 1--7.

\bibitem{mycielski-1955}
J.~Mycielski,
\emph{Sur le coloriage des graphes},
Colloq.\ Math.\ \textbf{3} (1955), 161--162.

\bibitem{sachs-1993}
H.~Sachs,
\emph{Elementary proof of the cycle-plus-triangles theorem},
Combinatorics, Paul Erd\H{o}s is Eighty, Bolyai Soc.\ Math.\ 
Stud., vol.~1, 1993, pp.~347--359.

\bibitem{santos-grt-2026}
V.\,F.\,S.~Santos,
\emph{Golden resolvent theory: spectral factorisation, Galois orbits, 
and Chebyshev ladders over $\Q(\sqrt{5})$},
Preprint, February 2026.

\bibitem{santos-golden-2026}
V.\,F.\,S.~Santos,
\emph{Golden eigenvalues and the Erd\H{o}s--Lov\'asz Tihany 
conjecture for Mycielski graphs},
Preprint, February 2026.

\bibitem{stiebitz-1987}
M.~Stiebitz,
\emph{Proof of a conjecture of T.~Gallai concerning connectivity 
properties of colour-critical graphs},
Combinatorica \textbf{7} (1987), 303--304.

\end{thebibliography}

\end{document}
